\documentclass[11pt]{amsart}

\usepackage{amsmath,amssymb,amsthm}
\usepackage[margin=27mm]{geometry}
\usepackage[hidelinks]{hyperref}

\newtheorem{theorem}{Theorem}[section]
\newtheorem{lemma}[theorem]{Lemma}
\newtheorem{proposition}[theorem]{Proposition}
\newtheorem{corollary}[theorem]{Corollary}
\theoremstyle{definition}
\newtheorem{definition}[theorem]{Definition}
\theoremstyle{remark}
\newtheorem{remark}[theorem]{Remark}

\newcommand{\Z}{\mathbb Z}
\newcommand{\Q}{\mathbb Q}
\newcommand{\gam}{\gamma}
\newcommand{\del}{\delta}
\newcommand{\ww}{\varpi}
\newcommand{\Hom}{\operatorname{Hom}}
\newcommand{\coker}{\operatorname{coker}}
\newcommand{\ann}{\operatorname{ann}}
\newcommand{\tr}{\operatorname{tr}}
\newcommand{\wt}{\operatorname{wt}}
\newcommand{\coeff}[1]{[#1]\,}
\newcommand{\pair}[2]{\langle #1,#2\rangle}

\title{The fifth dimension subring in nilpotency class three}
\author{}
\date{}

\begin{document}
\maketitle

\section{The statement and the standing reduction}

For a Lie ring \(L\), let \(U(L)\) be its universal enveloping
algebra and let \(\ww(L)\) be the augmentation ideal.  We use
\[
 \del_n(L)=L\cap\ww(L)^n,\qquad
 \gam_1(L)=L,\qquad \gam_{n+1}(L)=[L,\gam_n(L)].
\]
Our objective is
\begin{equation}\label{eq:main}
                         \del_5(L)\subseteq\gam_4(L)
\end{equation}
for every Lie ring \(L\).

We take the following previously established reduction as the starting
point of this paper.

\medskip
\noindent
\textbf{Standing reduction.}
To prove \eqref{eq:main}, it is enough to prove
\(\del_5(L)=0\) for finite Lie rings \(L\) whose additive group is a
\(2\)-group and which satisfy
\[
 \gam_4(L)=0,\qquad
 \gam_3(L)=\langle\bar z\rangle\cong\Z/q,\qquad q=2^\kappa .
\tag{R}
\]
\medskip

Everything after this reduction is proved below.  We use two standard
published facts: \(\del_3(L)=\gam_3(L)\) for every Lie ring, and the
free-presentation criterion
\begin{equation}\label{eq:presentation-criterion}
 a\in\del_n(L)
 \quad\Longleftrightarrow\quad
 f\in\ww(F)^n+\mathfrak r
\end{equation}
when \(F\twoheadrightarrow L=F/\mathcal R\) is a free presentation,
\(f\in F\) lifts \(a\), and
\(\mathfrak r=U(F)\mathcal R U(F)\).  These are respectively
\cite[Theorem~4.5 and Proposition~4.2]{BP}.

The proof after \textup{(R)} has three parts.  First, PBW coefficient
collection associates to each element of \(\del_5(L)\) a short matrix
system.  Second, a certificate satisfying \textup{(P1)--(P4)}
annihilates that system.  Third, a finite symplectic extension
constructs such a certificate for every admissible datum.

\section{An adapted presentation}

Until the final section, \(L\) satisfies \textup{(R)}.  Choose
\(x_1,\ldots,x_m\in L\) whose images form a Smith basis
\[
 L/\gam_2(L)\cong\bigoplus_{i=1}^m\Z/d_i,
 \qquad d_1\mid\cdots\mid d_m,
\]
and choose lifts, denoted by the same letters.  The \(d_i\) are powers
of \(2\).  These elements generate \(L\): if \(H\) is the Lie subring
they generate, then \(L=H+\gam_2(L)\); successively taking brackets
gives \(\gam_2(L)\subseteq H+\gam_3(L)\) and
\(\gam_3(L)\subseteq H+\gam_4(L)=H\).

Let \(F\) be the free Lie ring on the \(x_i\), let
\(p:F\twoheadrightarrow L\), and put
\(\mathcal R=\ker p\).  Write
\[
                         F=\bigoplus_{r\geq1}F_r
\]
for the grading by bracket length.  Since \(\gam_4(L)=0\),
\(\bigoplus_{r\geq4}F_r=\gam_4(F)\subseteq\mathcal R\).

Choose a basis of \(F_2\) consisting of elements
\[
 y_1,\ldots,y_n,\ \widetilde y_1,\ldots,\widetilde y_s
\]
such that the images of the \(y_k\) form a Smith basis
\[
 \gam_2(L)/\gam_3(L)\cong\bigoplus_{k=1}^n\Z/e_k,
 \qquad e_1\mid\cdots\mid e_n,
\]
and the kernel of \(F_2\to\gam_2(L)/\gam_3(L)\) is generated by
\(e_ky_k\) and the \(\widetilde y_s\).  Choose a basis of \(F_3\)
consisting of \(z,\widetilde z_1,\ldots,\widetilde z_t\), where
\[
 p(z)=\bar z,\qquad
 \ker(F_3\to\gam_3(L))
   =\langle qz,\widetilde z_1,\ldots,\widetilde z_t\rangle_{\Z}.
\]
The Smith choices give integers \(b_{ik},c_i,m_k,\widetilde c_s\)
and the following elements of \(\mathcal R\):
\begin{align}
 r_i&=d_ix_i-B_i-c_iz,
       &B_i&=\sum_k b_{ik}y_k, \label{eq:ri}\\
 s_k&=e_ky_k-m_kz, &
 \rho_s&=\widetilde y_s-\widetilde c_sz, \label{eq:sk}\\
 t&=qz, &
 \tau_a&=\widetilde z_a. \label{eq:t}
\end{align}

\begin{lemma}[relations through weight four]\label{lem:relations}
Choose a basis \(Q_1,\ldots,Q_N\) of \(F_4\).  Every element of
\(\mathcal R\), modulo \(\gam_5(F)\), is an integral linear
combination of
\[
 r_i,\quad s_k,\quad\rho_s,\quad t,\quad\tau_a,\quad Q_\nu .
\]
Moreover, if the element lies in
\(\mathcal R\cap\gam_c(F)\), the combination may be chosen using only
relation representatives whose least homogeneous degree is at least
\(c\).
\end{lemma}

\begin{proof}
Let \(r\in\mathcal R\).  Its degree-one component maps to zero in
\(L/\gam_2(L)\), hence is a linear combination of the \(d_ix_i\).
Subtracting the same combination of the \(r_i\) removes degree one.
The remaining degree-two component lies in the kernel of
\(F_2\to\gam_2(L)/\gam_3(L)\); subtracting \(s_k,\rho_s\) removes it.
The degree-three component then lies in the kernel of
\(F_3\to\gam_3(L)\); subtracting \(t,\tau_a\) removes it.  What remains
lies in \(\gam_4(F)\), and modulo \(\gam_5(F)\) it is a combination of
the \(Q_\nu\).  If \(r\in\gam_c(F)\), its components in degrees below
\(c\) are already zero, so the corresponding earlier subtraction
steps are empty.  This proves the last assertion.
\end{proof}

Put
\[
 D=\operatorname{diag}(d_1,\ldots,d_m),\qquad
 E=\operatorname{diag}(e_1,\ldots,e_n),\qquad
 B=(b_{ik}).
\]
Expand in the chosen basis of \(F_3\):
\[
 [x_i,y_k]=g_{ik}z+\sum_a h_{ika}\widetilde z_a,
\]
and set \(G=(g_{ik})\) and
\[
                              \Gamma=BG^{\mathsf T}.
\]
All matrix congruences below are modulo \(q\).

\begin{lemma}[the data identities]\label{lem:data}
The matrices \(D,E,B,G\) satisfy
\begin{equation}\label{eq:data-identities}
 DG=0,\qquad GE=0,\qquad
 \Gamma_{ii}=0,\qquad
 \Gamma_{ji}+\frac{d_j}{d_i}\Gamma_{ij}=0\quad(i<j).
\end{equation}
\end{lemma}

\begin{proof}
In \(L\), the element \(\bar z\) is central and
\([\gam_2(L),\gam_2(L)]\subseteq\gam_4(L)=0\).  Therefore
\[
 d_i[\bar x_i,\bar y_k]=[\bar B_i+c_i\bar z,\bar y_k]=0,
 \qquad
 e_k[\bar x_i,\bar y_k]=[\bar x_i,m_k\bar z]=0,
\]
which gives \(DG=0\) and \(GE=0\).  Moreover,
\begin{equation}\label{eq:bracket-data}
 d_p[\bar x_q,\bar x_p]
 =[\bar x_q,\bar B_p]
 =\Gamma_{pq}\bar z.
\end{equation}
Taking \(p=q=i\) gives \(\Gamma_{ii}=0\).  If \(i<j\), applying
\eqref{eq:bracket-data} to \((p,q)=(j,i)\) and \((i,j)\), and using
\(d_j=(d_j/d_i)d_i\), gives the last identity in
\eqref{eq:data-identities}.
\end{proof}

\section{The coefficient system for an element of
\texorpdfstring{\(\del_5\)}{delta 5}}

We now prove the normal-form information that will be used later.
This section is self-contained; in particular, no unpublished normal
form is being cited.

Give a homogeneous Lie element of \(F_r\) weight \(r\).  Order the
chosen homogeneous bases by
\[
 x_1<\cdots<x_m<y_1<\cdots<y_n
 <\widetilde y_1<\cdots<z<\widetilde z_1<\cdots<Q_1<\cdots .
\]
The filtered PBW theorem \cite[Chapter~I, \S2.7]{Bou} says that the
ordered products of these elements of total weight \(<5\), together
with \(1\), form a \(\Z\)-basis of \(U(F)/\ww(F)^5\).  For such a PBW
monomial \(\mu\), write \(\coeff{\mu}\xi\) for its coefficient.

For a strictly upper-triangular matrix \(u=(u_{ij})\), put
\begin{equation}\label{eq:ubar}
 \overline u=u-Du^{\mathsf T}D^{-1},\qquad
 \overline u_{ij}=
 \begin{cases}
  u_{ij},&i<j,\\
  -(d_i/d_j)u_{ji},&i>j,\\
  0,&i=j.
 \end{cases}
\end{equation}
This is integral because \(d_j\mid d_i\) when \(j<i\).  For matrices
of the same size use the Frobenius pairing
\[
                         \pair{X}{Y}=\tr(X^{\mathsf T}Y).
\]

\begin{lemma}[PBW collection]\label{lem:collection}
Let \(f\in\gam_3(F)\) and \(\theta\in\mathfrak r\) satisfy
\(\theta-f\in\ww(F)^5\).  Suppose that the image of \(f\) in \(L\) is
\(\zeta\bar z\).  Then there exist integer matrices
\[
 u=(u_{ij})_{i<j},\qquad v,v'\in M_{m\times n}(\Z)
\]
such that
\begin{align}
 \overline uB+Dv+v'E&=0, \tag{B}\label{eq:B}\\
 \sum_{i<j}u_{ij}\frac{d_j}{d_i}\Gamma_{ij}
   &\equiv\zeta\pmod q, \tag{Z}\label{eq:Z}\\
 (v^{\mathsf T}B+B^{\mathsf T}v)_{kl}
   &\equiv0\pmod{\gcd(e_k,e_l)}\quad(k<l), \tag{C1}\label{eq:C1}\\
 (v^{\mathsf T}B)_{kk}
   &\equiv0\pmod{e_k}. \tag{C2}\label{eq:C2}
\end{align}
\end{lemma}

\begin{proof}
We first explain completely how relation-products are collected.
By Lemma~\ref{lem:relations}, modulo \(\ww(F)^5\), every element of
\(\mathfrak r\) is a sum of products
\begin{equation}\label{eq:relation-product}
             a\,\sigma\,b,\qquad
 \sigma\in\{r_i,s_k,\rho_s,t,\tau_a,Q_\nu\},
\end{equation}
where \(a,b\) are PBW products and only the following external weights
can occur:
\begin{center}
\begin{tabular}{c|c|l}
\(\sigma\) & its least weight & possible total external weight\\ \hline
\(r_i\) & \(1\) & \(0,1,2,3\)\\
\(s_k,\rho_s\) & \(2\) & \(0,1,2\)\\
\(t,\tau_a\) & \(3\) & \(0,1\)\\
\(Q_\nu\) & \(4\) & \(0\).
\end{tabular}
\end{center}
This table is exhaustive in total weight at most four.

There is no hidden commutation step in this assertion.  If \(h\) is
one homogeneous Lie-basis element and \(\sigma\in\mathcal R\), then
\begin{equation}\label{eq:move-relation}
                    h\sigma=\sigma h+[h,\sigma],
 \qquad [h,\sigma]\in\mathcal R.
\end{equation}
Thus a relation can be moved through a PBW product one Lie factor at
a time.  The new relation in \eqref{eq:move-relation} has at least
the sum of the two weights and one fewer external factor; the final assertion of
Lemma~\ref{lem:relations} decomposes it without lowering that weight.
Likewise, an inverted pair of ordinary
Lie factors is ordered by
\begin{equation}\label{eq:pbw-swap}
                    h'h=hh'-[h,h']\qquad(h<h'),
\end{equation}
and the commutator has fewer PBW factors.  Induction first on total
external factors and then on the number of inversions therefore puts
the least-weight PBW word in every term of
\eqref{eq:relation-product} into ordered form:
the reordered term has one fewer inversion, while the commutator term
has one fewer external factor.  The process terminates, and every
term created by it is again in the displayed table.

Fix one fully collected representation and combine equal canonical
products.  We now read it weight by weight.  In weight one, only a scalar
multiple of some \(r_i\) can occur.  Its coefficient at \(x_i\) is a
multiple of \(d_i\).  Since \(f\in\gam_3(F)\), all such coefficients
are zero; hence no \(r_i\)-term remains.

In weight two, the only two-factor words come from \(r_i\) with one
external \(x\).  For \(i<j\), their ordered representatives are
\(r_ix_j\) and \(x_ir_j\); for \(i=j\) use \(r_ix_i\).  The diagonal
coefficient is a multiple of \(d_i\), so it is zero.  For \(i<j\),
if the two coefficients are \(a_{ij}\) and \(c_{ij}\), the coefficient
of \(x_ix_j\) is
\[
                         d_i a_{ij}+d_jc_{ij}=0.
\]
Since \(d_i\mid d_j\), this pair is an integral multiple of
\[
                    \frac{d_j}{d_i}r_ix_j-x_ir_j.
\]
The remaining one-factor weight-two terms are multiples of \(s_k\)
and \(\rho_s\); their leading terms \(e_ky_k\) and
\(\widetilde y_s\) are independent, so their coefficients are zero.
Consequently the entire part whose least weight is below three is a
linear combination
\begin{equation}\label{eq:ThetaII}
 \Theta_{II}=
 \sum_{i<j}u_{ij}
 \left(\frac{d_j}{d_i}r_ix_j-x_ir_j\right).
\end{equation}
Every summand in \eqref{eq:ThetaII} has weight at least three.

At weight three, the canonical relation-products, apart from
\eqref{eq:ThetaII}, are
\begin{equation}\label{eq:weight-three-list}
 r_i y_k,\quad x_i s_k,
\quad r_i\widetilde y_s,\quad x_i\rho_s,
\quad t,\quad\tau_a.
\end{equation}
and the following three placements:
\[
 r_i x_jx_k\ (i\leq j\leq k),\qquad
 x_i r_jx_k\ (i<j\leq k),\qquad
 x_ix_jr_k\ (i\leq j<k).
\]
They are precisely the placements for which the leading
three-\(x\) word is ordered; cases of equal adjacent indices are
assigned to the family on their left.  Every other placement reduces
to these by \eqref{eq:move-relation}--\eqref{eq:pbw-swap}, and its
correction terms occur in \eqref{eq:weight-three-list}, in the three
displayed families, or first in weight four.
Let \(v_{ik}\) and
\(v'_{ik}\) be the coefficients of \(r_iy_k\) and \(x_is_k\), and put
\[
                  \Theta=\Theta_{II}
                    +\sum_{i,k}v_{ik}r_iy_k
                    +\sum_{i,k}v'_{ik}x_is_k.
\]

We record the three PBW coefficients relevant below.  Directly from
\eqref{eq:ri}--\eqref{eq:sk}, reordering
\(y_kx_j=x_jy_k-[x_j,y_k]\), one obtains
\begin{align}
 \coeff{x_i y_k}\Theta
   &=(\overline uB+Dv+v'E)_{ik},\label{eq:coeff-xy}\\
 \coeff{z}\Theta
   &=\sum_{i<j}u_{ij}\frac{d_j}{d_i}\Gamma_{ij},
       \label{eq:coeff-z}\\
 \coeff{y_ky_l}\Theta
   &=-(v^{\mathsf T}B+B^{\mathsf T}v)_{kl}
       &&(k<l),\label{eq:coeff-yy}\\
 \coeff{y_k^2}\Theta
   &=-(v^{\mathsf T}B)_{kk}.\label{eq:coeff-square}
\end{align}
For clarity, in the summand indexed by \(i<j\) in
\eqref{eq:ThetaII},
\(-\frac{d_j}{d_i}B_ix_j\) contributes
\(-\frac{d_j}{d_i}b_{ik}\) at \(x_jy_k\) and
\(\frac{d_j}{d_i}\Gamma_{ij}\) at \(z\), while
\(x_iB_j\) contributes \(b_{jk}\) at \(x_iy_k\).
This proves both the signs and the transpose in
\eqref{eq:coeff-xy}--\eqref{eq:coeff-z}.  Equations
\eqref{eq:coeff-yy}--\eqref{eq:coeff-square} come only from
\(-\sum_{i,k}v_{ik}B_iy_k\).

No other weight-three term contributes to \(x_iy_k\).  Indeed,
\(r_i\widetilde y_s\) and \(x_i\rho_s\) have only
\(x_i\widetilde y_s\) as a two-factor term in weight three.  Any
correction created while moving a supplementary relation lies in
\(\mathcal R\cap\gam_3(F)\), so Lemma~\ref{lem:relations} shows that
its named \(z\)-coefficient is a multiple of \(q\).  Each of the three
ordered \(r\)-and-two-\(x\) families has a three-\(x\) leading word,
and its \(B_i\)-part first occurs in weight four.  Ordering that part
produces only \(xxY\), \(xZ\), or a single degree-four Lie factor,
never a pure \(yy\)-product.  Finally,
\(t,\tau_a\) are single Lie factors.  Since a Lie element has no
multi-factor PBW
coefficient, \(\coeff{x_iy_k}f=0\).  Comparing with
\(\theta-f\in\ww^5\) gives \eqref{eq:B}.

Likewise, among the remaining terms of weight three, only a multiple
of \(t=qz\) can change the coefficient of \(z\); \(\tau_a\) changes
only a supplementary coefficient.  The \(z\)-coefficient of \(f\)
is congruent to \(\zeta\) modulo \(q\), because the supplementary
degree-three basis elements and \(F_{\geq4}\) vanish in \(L\).
Equation \eqref{eq:coeff-z} therefore gives \eqref{eq:Z}.

It remains to justify that no omitted weight-four term spoils
\eqref{eq:coeff-yy}--\eqref{eq:coeff-square}.  After the
factor-by-factor collection
\eqref{eq:move-relation}--\eqref{eq:pbw-swap}, the exhaustive shape
table is as follows.  Here
\(Y\) denotes a named or supplementary degree-two basis element and
\(Z\) a degree-three basis element; ``pure \(yy\)'' means a product
of two of the named \(y_k\)'s.
\begin{equation}\label{eq:shape-table}
\begin{array}{c|c|c|c}
\text{relation}&\text{external factors}&
 \text{ordered leading shape}&\text{pure \(yy\) part}\\ \hline
r_i&xxx&xxxx&0\\
r_i&x,Y&xxY&0\\
r_i&Z&xZ&0\\
s_k,\rho_s&xx&xxY&0\\
s_k&\widetilde y_a&y_k\widetilde y_a&0\\
\rho_s&Y&\widetilde y_sY&0\\
t,\tau_a&x&xZ&0\\
Q_\nu&1&Q_\nu&0.
\end{array}
\end{equation}
The missing row is \(s_k\) with one named \(y\)-factor.  Its ordered
forms are exactly
\begin{equation}\label{eq:s-y-products}
 s_ky_l=e_ky_ky_l+\text{(terms with \(z\))},\qquad
 y_ks_l=e_ly_ky_l+\text{(terms with \(z\))}\quad(k<l).
\end{equation}
and \(s_ky_k=e_ky_k^2+\text{(terms with \(z\))}\).

This also checks every correction created by collection.  At total
weight four, a commutator term in \eqref{eq:move-relation} has fewer
external factors.
If its new relation has least weight two, its remaining external
weight is two; Lemma~\ref{lem:relations} makes it a combination of
\(s_k,\rho_s\), followed either by two \(x\)'s or by one
degree-two factor.  These are the fourth and fifth rows of
\eqref{eq:shape-table}, and the terms in \eqref{eq:s-y-products}.
If the new relation has least weight three, it has one external
\(x\) and is a combination of \(t,\tau_a\); if it has weight four,
it is a Lie element \(Q_\nu\).  A correction in
\eqref{eq:pbw-swap} replaces two
external \(x\)'s by one degree-two factor, or an external \(x,Y\)
pair by one degree-three factor.  It therefore stays in the
corresponding \(r_i\)-row, or takes an \(s_k,\rho_s\)-with-two-\(x\)
term to its \(s_kY,\rho_sY\) row.  Thus the table is closed under both
\eqref{eq:move-relation} and \eqref{eq:pbw-swap}.  The potentially
dangerous reduction of
\(r_i\) with three \(x\)'s is included here: replacing two inverted
\(x\)'s by their degree-two bracket leaves either an ordered
\(r_i xY\)-term, whose leading shape is \(xxY\), or a term
\([x,r_i]Y\); the latter has a weight-two relation and is therefore a
combination of the \(s_kY,\rho_sY\) rows.
In particular, every omitted contribution to
\(\coeff{y_ky_l}\theta\) lies in
\[
 e_k\Z+e_l\Z=\gcd(e_k,e_l)\Z,
\]
and every omitted contribution to \(\coeff{y_k^2}\theta\) lies in
\(e_k\Z\): an \(s_aY\)-term can meet the named product \(y_ky_l\)
only when \(a=k\) or \(a=l\), and on the diagonal only when \(a=k\).

Finally, the degree-four part of the Lie element \(f\) is a
combination of the single Lie-basis factors \(Q_\nu\); hence its
coefficients at \(y_ky_l\) and \(y_k^2\) are zero.  Comparing
\eqref{eq:coeff-yy}--\eqref{eq:s-y-products} with
\(\theta-f\in\ww^5\) proves
\eqref{eq:C1}--\eqref{eq:C2}.
\end{proof}

\begin{theorem}[coefficient system]\label{thm:system}
If \(a=\zeta\bar z\in\del_5(L)\), then there exist \(u,v,v'\)
satisfying \eqref{eq:B}--\eqref{eq:C2}.
\end{theorem}

\begin{proof}
Since \(\del_5(L)\subseteq\del_3(L)=\gam_3(L)\), choose
\(f\in\gam_3(F)\) mapping to \(a\).  By
\eqref{eq:presentation-criterion}, there is
\(\theta\in\mathfrak r\) with \(\theta-f\in\ww(F)^5\).
Apply Lemma~\ref{lem:collection}.
\end{proof}

\section{Certificates annihilate
\texorpdfstring{\(\del_5\)}{delta 5}}

\begin{definition}\label{def:certificate}
A \emph{certificate} for \(D,E,B,G\) is a pair
\[
 \alpha\in M_{m\times n}(\Z/q),\qquad
 \Lambda=\Lambda^{\mathsf T}\in M_n(\Z/q)
\]
such that, with \(X=B\alpha^{\mathsf T}\), one has
\begin{align}
 X_{ji}-\frac{d_j}{d_i}X_{ij}&=\Gamma_{ji}
                                   &&(i<j), \tag{P1}\label{eq:P1}\\
 D\alpha&=B\Lambda,                 &&       \tag{P2}\label{eq:P2}\\
 \alpha E&=0,                       &&       \tag{P3}\label{eq:P3}\\
 E\Lambda&=0.                       &&       \tag{P4}\label{eq:P4}
\end{align}
\end{definition}

\begin{theorem}[certificate criterion]\label{thm:criterion}
If \(D,E,B,G\) admit a certificate, then \(\del_5(L)=0\).
\end{theorem}

\begin{proof}
Let \(a=\zeta\bar z\in\del_5(L)\), and take \(u,v,v'\) from
Theorem~\ref{thm:system}.  Put
\[
 N=\sum_{i<j}u_{ij}\Gamma_{ji}
   =\pair{u}{\Gamma^{\mathsf T}}.
\]
By \eqref{eq:Z} and Lemma~\ref{lem:data},
\begin{equation}\label{eq:zetaN}
                            \zeta\equiv-N\pmod q.
\end{equation}
Since
\[
 X_{ji}=(\alpha B^{\mathsf T})_{ij},
 \qquad X_{ij}=(\alpha B^{\mathsf T})_{ji},
\]
condition \eqref{eq:P1} and the definition of \(\overline u\) give
\[
\begin{aligned}
 N
 &=\sum_{i<j}u_{ij}
   \left((\alpha B^{\mathsf T})_{ij}
       -\frac{d_j}{d_i}(\alpha B^{\mathsf T})_{ji}\right)\\
 &=\pair{\overline u}{\alpha B^{\mathsf T}}
  =\pair{\overline uB}{\alpha}.
\end{aligned}
\]
Using \eqref{eq:B} and then \eqref{eq:P2}--\eqref{eq:P3},
\[
 N=-\pair{Dv+v'E}{\alpha}
   =-\pair{v}{D\alpha}-\pair{v'}{\alpha E}
   =-\pair{v}{B\Lambda}.
\]
Together with \eqref{eq:zetaN}, this yields
\begin{equation}\label{eq:zeta-lambda}
                   \zeta\equiv\pair{\Lambda}{B^{\mathsf T}v}.
\end{equation}

Because \(\Lambda\) is symmetric,
\begin{equation}\label{eq:lambda-expansion}
 \pair{\Lambda}{B^{\mathsf T}v}
 =\sum_k\Lambda_{kk}(v^{\mathsf T}B)_{kk}
  +\sum_{k<l}\Lambda_{kl}
       (v^{\mathsf T}B+B^{\mathsf T}v)_{kl}.
\end{equation}
Condition \(E\Lambda=0\), together with symmetry, gives
\[
 e_k\Lambda_{kl}=e_l\Lambda_{kl}=0\quad\text{in }\Z/q.
\]
B\'ezout's identity therefore gives
\(\gcd(e_k,e_l)\Lambda_{kl}=0\).  Every summand in
\eqref{eq:lambda-expansion}
vanishes by \eqref{eq:C1}--\eqref{eq:C2}.  Hence
\(\zeta=0\) in \(\Z/q\), so \(a=0\).
\end{proof}

\section{Existence of certificates}

We now prove that every datum satisfying
\eqref{eq:data-identities} has a certificate.  The proof is invariant
until its final matrix translation.

Put \(C=\Z/q\), embedded in \(\Q/\Z\) by
\(\bar r\mapsto r/q+\Z\).  If \(A\) is a finite abelian group of
exponent dividing \(q\), write \(A^\vee=\Hom(A,C)\).  Evaluation is a
canonical isomorphism \(A\cong A^{\vee\vee}\), as is checked on each
cyclic summand \(\Z/d\) with \(d\mid q\).
An \emph{integer lift} of a \(C\)-valued bilinear form is an
integer-valued bilinear form whose reduction modulo \(q\) is the given
form.

\begin{lemma}[polarization]\label{lem:polarization}
Every alternating bilinear form
\(\Omega:\mathcal E\times\mathcal E\to\Q/\Z\) on a finite abelian
group is of the form
\[
                              H^{\mathsf T}-H=\Omega
\]
for some bilinear \(H:\mathcal E\times\mathcal E\to\Q/\Z\).
\end{lemma}

\begin{proof}
Choose a cyclic decomposition
\(\mathcal E=\bigoplus_{r=1}^s\langle e_r\rangle\).  For \(r<t\), set
\[
 H(e_r,e_t)=0,\qquad
 H(e_t,e_r)=\Omega(e_r,e_t),\qquad H(e_r,e_r)=0,
\]
and extend bilinearly.  This is well defined because the orders of
both generators annihilate \(\Omega(e_r,e_t)\), and the displayed
identity follows on generators.  No division by \(2\) is used.
\end{proof}

\begin{theorem}[abstract certificate]\label{thm:abstract}
Let \(P\) be a finite-rank free abelian group, let
\(d:P\to P\) be injective with finite cokernel, and let
\[
                    P\xrightarrow{d}P\xrightarrow{\pi}A
\]
satisfy \(\pi d=0\).  Let \(b:P\to A^\vee\), and define
\[
                         \varphi(x,z)=b(x)(\pi z).
\]
Suppose that \(\varphi\) has an integer lift
\(\widehat\varphi:P\times P\to\Z\) satisfying
\begin{equation}\label{eq:lift-skew}
 \widehat\varphi(x,dy)+\widehat\varphi(y,dx)=0
 \qquad(x,y\in P).
\end{equation}
Then there exist a homomorphism \(\alpha:P\to A\), a symmetric
bilinear form \(\lambda:A^\vee\times A^\vee\to C\), and an integer
lift \(\widehat\psi\) of
\(\psi(x,z)=b(x)(\alpha z)\), such that
\begin{align}
 \widehat\psi(x,dy)-\widehat\psi(y,dx)
   &=\widehat\varphi(x,dy), \label{eq:abstract-psi}\\
 \alpha d&=\lambda_*b.     \label{eq:abstract-alpha}
\end{align}
Here \(\lambda_*(\eta)\in A=A^{\vee\vee}\) is defined by
\(\chi(\lambda_*(\eta))=\lambda(\eta,\chi)\).
\end{theorem}

\begin{proof}
Extend \(d\) and \(\widehat\varphi\) to \(P\otimes\Q\).  Since \(d\)
is invertible over \(\Q\), the rule
 \begin{equation}\label{eq:vartheta}
                     \vartheta(dx,z)
                        =\frac{\widehat\varphi(x,z)}q
 \end{equation}
defines a rational bilinear form on \(P\otimes\Q\).
Equation \eqref{eq:lift-skew} makes it skew on \(dP\times dP\);
because \(dP\) spans \(P\otimes\Q\), \(\vartheta\) is alternating.

Form the finite abelian group
\[
 \mathcal E=
 \frac{A^\vee\oplus P}{\{(b(x),-dx):x\in P\}},
 \qquad i(\chi)=[\chi,0],\qquad j(z)=[0,z].
\]
The sequence
\[
 0\longrightarrow A^\vee\xrightarrow{i}\mathcal E
   \longrightarrow\coker d\longrightarrow0
\]
is exact: injectivity of \(i\) follows from injectivity of \(d\), and
the map \([\chi,z]\mapsto z+dP\) has kernel \(i(A^\vee)\).
Also \(jd=ib\).  On \(A^\vee\oplus P\), define
\[
 \Omega_0((\chi,z),(\eta,w))
  =\chi(\pi w)-\eta(\pi z)+\vartheta(z,w)\pmod{\Z}.
\]
For a defining relation,
\[
\begin{aligned}
 \Omega_0((b(x),-dx),(\eta,w))
 &=b(x)(\pi w)+\eta(\pi dx)
      -\frac{\widehat\varphi(x,w)}q\\
 &=0\quad\text{in }\Q/\Z.
\end{aligned}
\]
Here \(\pi d=0\), and the first and last terms agree because
\(\widehat\varphi\) lifts \(\varphi\).  The form \(\Omega_0\) is
alternating, so the relations vanish in both variables.  Thus it
descends to an alternating form \(\Omega\) on \(\mathcal E\).

By Lemma~\ref{lem:polarization}, choose \(H\) with
\(H^{\mathsf T}-H=\Omega\).  Define \(\alpha\) and \(\lambda\) by
\begin{equation}\label{eq:alpha-lambda}
 \chi(\alpha z)=H(jz,i\chi),\qquad
 \lambda(\chi,\eta)=H(i\eta,i\chi).
\end{equation}
The values lie in
\((\Q/\Z)[q]=C\), because every \(i\chi\) is killed by \(q\).
Hence the first formula indeed defines an element of
\(A=A^{\vee\vee}\).  Since \(\Omega\) vanishes on
\(i(A^\vee)\times i(A^\vee)\), the second formula is symmetric.
Moreover,
\[
\begin{aligned}
 \chi(\alpha dy)
 &=H(jdy,i\chi)=H(ib(y),i\chi)\\
 &=H(i\chi,ib(y))
  =\lambda(b(y),\chi)
  =\chi(\lambda_*b(y)),
\end{aligned}
\]
which proves \eqref{eq:abstract-alpha}.

Choose a rational bilinear lift
\[
 S:P\times P\to\Q,\qquad
 S(z,w)\pmod{\Z}=H(jz,jw).
\]
Then
\[
                       S^{\mathsf T}-S\equiv\vartheta\pmod\Z.
\]
The difference \(\Delta=S^{\mathsf T}-S-\vartheta\) is an integral
alternating form.  Choose an integral triangular matrix \(K\) with
\(K^{\mathsf T}-K=-\Delta\), which is possible by putting the entries
of \(-\Delta\) into one triangle.  After replacing \(S\) by \(S+K\),
we may therefore suppose that
\begin{equation}\label{eq:exact-S}
                           S^{\mathsf T}-S=\vartheta
\end{equation}
exactly over \(\Q\).  Set
\[
                         \widehat\psi(x,z)=qS(z,dx).
\]
This is integral because, modulo \(\Z\),
\[
 S(z,dx)=H(jz,jdx)=H(jz,ib(x))=b(x)(\alpha z)\in C.
\]
The same equality says that \(\widehat\psi\) lifts \(\psi\).  Finally,
by \eqref{eq:vartheta} and \eqref{eq:exact-S},
\[
\begin{aligned}
 \widehat\psi(x,dy)-\widehat\psi(y,dx)
 &=q\bigl(S(dy,dx)-S(dx,dy)\bigr)\\
 &=q\vartheta(dx,dy)
  =\widehat\varphi(x,dy),
\end{aligned}
\]
which proves \eqref{eq:abstract-psi}.
\end{proof}

\begin{theorem}[matrix certificate]\label{thm:matrix}
Let \(D,E,B,G\) be as in Section~2 and suppose only that they satisfy
\eqref{eq:data-identities}.  Then they admit a certificate in the
sense of Definition~\ref{def:certificate}.
\end{theorem}

\begin{proof}
Let \(R_q=\Z/q\), use row vectors, and put
\[
                       A=\{a\in R_q^{1\times n}:aE=0\}.
\]
The dot product induces a perfect pairing
\begin{equation}\label{eq:pairing}
 \frac{R_q^{1\times n}}{R_q^{1\times n}E}\times A\longrightarrow R_q,
 \qquad(\bar v,a)\longmapsto va^{\mathsf T}.
\end{equation}
Indeed, in one coordinate, if \(g=\gcd(e,q)\), then
\[
 \ann_{R_q}(e)=(q/g)R_q,\qquad
 \ann_{R_q}(\ann_{R_q}(e))=gR_q=eR_q.
\]
Applying this coordinatewise proves that the left kernel in
\eqref{eq:pairing} is exactly \(R_q^{1\times n}E\), and the resulting
finite groups have the same order: coordinatewise both orders equal
\(\gcd(e,q)\), and \(|A^\vee|=|A|\).  Thus the first factor is
\(A^\vee\).

Take \(P=\Z^m\), with basis \(x_1,\ldots,x_m\), and let \(d\) have
matrix \(D\).  Define
\[
                    \pi(x_i)=G_i\in A,\qquad
                    b(x_i)=\overline{B_i}\in A^\vee,
\]
where \(G_i,B_i\) are rows.  The identities \(GE=0\) and \(DG=0\)
give respectively \(\pi(x_i)\in A\) and \(\pi d=0\).
The form \(b_\pi\) has matrix \(BG^{\mathsf T}=\Gamma\).

For \(i<j\), choose an integer \(r_{ij}\) representing
\(\Gamma_{ij}\), and define
\begin{equation}\label{eq:hatphi}
 \widehat\varphi_{ij}=r_{ij},\qquad
 \widehat\varphi_{ji}=-\frac{d_j}{d_i}r_{ij},\qquad
 \widehat\varphi_{ii}=0.
\end{equation}
The last two identities in \eqref{eq:data-identities} say that this
is an integer lift of \(\Gamma\), and
\[
 d_j\widehat\varphi_{ij}+d_i\widehat\varphi_{ji}=0.
\]
Theorem~\ref{thm:abstract} applies.

Let the row \(\alpha_i\) represent \(\alpha(x_i)\in A\), and define
\[
 \Lambda_{kl}=\lambda(\bar e_k,\bar e_l),
\]
where the \(\bar e_k\) are the standard generators of the quotient
in \eqref{eq:pairing}.  Then
\[
 \alpha E=0,\qquad
 \Lambda=\Lambda^{\mathsf T},\qquad E\Lambda=0.
\]
Under \eqref{eq:pairing}, \(\lambda_*(\bar v)\) is represented by
the row \(v\Lambda\).  Hence \(\alpha d=\lambda_*b\) is precisely
\[
                               D\alpha=B\Lambda.
\]

Finally, \(\psi=b_\alpha\) has matrix \(X=B\alpha^{\mathsf T}\).
On \(x_i,x_j\), the exact equality \eqref{eq:abstract-psi} gives
\[
 d_j\widehat\psi_{ij}-d_i\widehat\psi_{ji}
   =d_j\widehat\varphi_{ij}.
\]
For \(i<j\), division by \(d_i\) and \eqref{eq:hatphi} give
\[
 \widehat\psi_{ji}-\frac{d_j}{d_i}\widehat\psi_{ij}
   =\widehat\varphi_{ji}.
\]
Reducing modulo \(q\) is exactly \eqref{eq:P1}.  Choosing symmetric
integer representatives for \(\Lambda\) and arbitrary representatives
for \(\alpha\) gives the required integer matrices if desired.
\end{proof}

\section{Conclusion}

\begin{theorem}\label{thm:minimal}
Every finite Lie ring satisfying \textup{(R)} has
\(\del_5(L)=0\).
\end{theorem}

\begin{proof}
Construct \(D,E,B,G\) as in Section~2.  Lemma~\ref{lem:data} gives
the hypotheses of Theorem~\ref{thm:matrix}, so a certificate exists.
Theorem~\ref{thm:criterion} then gives \(\del_5(L)=0\).
\end{proof}

\begin{corollary}[main result, after the standing reduction]
For every Lie ring \(L\),
\[
                              \del_5(L)\subseteq\gam_4(L).
\]
\end{corollary}

\begin{proof}
Apply Theorem~\ref{thm:minimal} to the reduced case \textup{(R)}, and
then apply the standing reduction from Section~1.
\end{proof}

\begin{thebibliography}{9}

\bibitem{BP}
L.~Bartholdi and I.~B.~S. Passi,
\emph{Lie dimension subrings},
J. Algebra Appl. \textbf{14} (2015), no.~6, 1550042.

\bibitem{Bou}
N.~Bourbaki,
\emph{Lie Groups and Lie Algebras, Chapters 1--3},
Springer, 1989.

\end{thebibliography}

\end{document}
